\section{Ringkasan Uji Monotonisitas (Kalibrasi)}

Uji monotonisitas dipakai untuk memverifikasi bahwa \textit{baseline} diimplementasikan secara sahih (kualitas tidak menurun ketika anggaran dinaikkan) sekaligus mengidentifikasi tuas hiperparameter yang benar-benar menggerakkan cakupan. 
Dua tuas diuji seragam pada seluruh algoritma.

\begin{longtable}{
    @{}
    >{\raggedright\arraybackslash}p{3cm}          % Tuas
    >{\raggedleft\arraybackslash}p{2.5cm}           % Nilai diuji
    >{\raggedright\arraybackslash}p{3.5cm}          % Efek terhadap CoDA-EDA
    >{\raggedright\arraybackslash}p{3.5cm}          % Kesimpulan
    @{}
}
\caption{Efek tuas hiperparameter terhadap kualitas (dataset e-document)} \\
\toprule
Tuas & Nilai diuji & Efek terhadap CoDA-EDA & Kesimpulan \\
\midrule
\endfirsthead
\caption{Efek tuas hiperparameter terhadap kualitas (dataset e-document) (lanjutan)} \\
Tuas & Nilai diuji & Efek terhadap CoDA-EDA & Kesimpulan \\
\midrule
\endhead
\bottomrule
\endfoot

Bobot false positive ($w_{\mathrm{fp}}$) & $1,0 / 2,0 / 4,0$ & Cakupan $0,186 \rightarrow 0,187$ (praktis datar) & Bukan tuas cakupan; ditetapkan $w_{\mathrm{fp}} = 1,0$ \\

Anggaran jumlah aturan & $4 / 8 / 12$ & F-score $0,362 \rightarrow 0,571 \rightarrow 0,568$ (monoton naik) & Tuas cakupan sesungguhnya; jumlah aturan diperlakukan sebagai parameter rekayasa \\

\end{longtable}

Karena kenaikan anggaran aturan menaikkan F-score secara monoton (tanpa penurunan), implementasi kerangka penambangan dinyatakan sahih. 
Nilai baku eksperimen ditetapkan $w_{\mathrm{fp}} = 1,0$; jumlah aturan efektif pada anggaran baku berkisar 4--5.